\documentclass{article}
\usepackage{amsmath}
\title{Equal cave relay sets do not determine equal native receive effects}
\author{}
\date{}
\begin{document}
\maketitle

\paragraph{Audit.}
At the same modeled cave junction, two deployment histories have
the same two static breadcrumbs, the same endpoints, and the same
minimum-cost relay tile, yet the native visibility model and broker
deliver a report in only one. A fresh model first receiving both
positions can disagree with a live model that learned the positions
sequentially. This is an I2 development comparison of a source
interface, not a field experiment, independently chosen mission, or
full stack test.

\paragraph{Source seam.}
The pinned SubT \texttt{UpdateIfNewBreadcrumbs} scans the observed
poses for previously unseen static names. When a new name appears,
it passes the complete accumulated position set to the visibility
model. The actual route-table method appends these positions to the
matching tile's breadcrumb list without replacing its first entry.
Route construction then uses that entry for hop geometry.
In the original simulation this method concerns named static
models, so the experiment uses additions only. Neither pose edits
to an installed model, stock revocation, nor live failure detection
is used as an automatic removal mechanism.

The replay imports the exact published Cave Circuit 01 visibility
graph and SDF tile origins. It computes all shortest visibility
costs on the 158-vertex graph with Dijkstra and an independent
Floyd--Warshall recurrence, then passes their exact three-tile
projection to the original route update. The exact converter from
world coordinates to a tile is not imported: the four distinct
sample points are explicitly mapped to the stated tiles in a
declared four-point lookup seam. The sample positions one metre
from the tile origin are development probes, not surveyed
traversability, collider clearance, nor verified coordinates in a
production LUT. For the source witness there are only two proposed
static breadcrumbs and both are in the same projected tile, so
omitting other vertices from the relay enumeration does not remove
any allowed relay. The full graph has already contributed to every
visibility cost in the projection.

The source extracts also run the original route-to-power calculation,
transmission and receiver accounting, broker dispatch, and bound
receive callback. The selected RF envelope has high transmit power,
zero fading variance, fixed admissible RNG draw, and an explicitly
set hard range. No measured packet probabilities are inferred from
this envelope. No hard range used here is attributed to field
hardware. Separate traffic is dispatched on its own unaffected
native route in every test.

\paragraph{Witness and exact boundary.}
Use tiles $3,2,39$, with visibility costs $12$ direct and $6,6$
through tile~2. At the declared visibility-cost cap of~10, no
direct transmission is possible. The endpoint origins are
$(75,50,0)$ and $(50,25,0)$. The two eligible positions in tile~2
are $a=(49,50,0)$ and $b=(51,50,0)$. The required geometric max hop
is $26$ via $a$ and $\sqrt{626}$ via $b$. For simulator range
parameter $R$ in the interval $[\sqrt{626},26)$, only $b$ qualifies:
the pinned code rejects a range strictly larger than the cap. At
$R\approx25.509996$ this is not a floating-point tie. Install $a$
first and the old first entry remains $a$ even after $b$ is present.
Reversing discovery selects $b$, and the same two-byte report reaches
its bound callback. A fresh table seeing both at once selects $a$
because the executed comparator chooses the smaller coordinate
when the two points differ on just that axis. The pinned
Gazebo Math comparator is not assumed to be lexicographic. Sole-relay and no-relay controls discriminate selection
from absolute unreachability. This is replacement by cold
reconstruction of a development state, not a command the stock
controller can necessarily issue during the mission.

The route's tile and graph-cost outputs have not changed between the
two valid complete discovery histories. The missing distinction is
the effective first representative per tile. Keeping the entire
ordered insertion trace is unnecessary on this append-only,
frozen-graph seam with consistently ordered candidate sets: preserve the original named position set for
detecting actual new registrations and the first representative
of each already populated tile. A new distinct name can add a tile
or supply its first representative, but it cannot change an existing
one. With these data and exact graph and lookup, reconstruction of
the route-cost run has the same sorted sets and reads the same
first entries; by induction each subsequent admitted addition
maintains those facts. This argument excludes graph/LUT updates,
per-name movement and deletion, position sets outside a strict weak ordering regime,
and non-static hardware behavior. A same-information direct
calculation ties; the comparison against aggregate replay shows lost
state, not a better route algorithm.

\paragraph{Verification.}
The program exhausts the stated imports and one-metre axis offsets,
keeping nondegenerate same-tile pairs. For each, it tests an interval
point between the two hop lengths and controls below and above,
and compares eight deployment and reconstruction shapes. In
particular, restarting with the two live positions alone is distinct
from seeding a cold table from the persisted first representative.
Both first-entry reconstructions reproduce their respective source
delivery results in every sampled cut. Results separate a selected
relay, finite modeled power, actual application delivery, and receive
debit. A selected relay can still fail the range condition. The
package supplies source hashes and a clean extraction replay.
\end{document}
